% ============================================================================
% Held-out validation of the calibrated networks (draft for the paper).
% Numbers from scripts/09_network_validation.R (2026-11-05); table generated by
% that script: output/09_network_validation/validation_table.tex; figure:
% plots/09_network_validation/network_validation.pdf.
% Bibliography keys used: mcpherson2001, smith2014, granovetter1973 (all in
% paper-NDFU/references.bib). To add before citing: Boase et al. (2006) for the
% mean degree of 30; Feld (1981) if the foci sentence is kept.
% ============================================================================
\subsection{Held-out validation of the networks}
\label{sec:validation}

The networks are built from two statistics of the GSS 2004 network module per racial
group: the mean number of confidants (relative degree) and the share of confidants of
the ego's own race. We ask whether they reproduce features of real ego networks that
were not used to build them. The 1{,}424 respondents of the module with race and number
of confidants known were split into two halves, stratified by race. Half A supplied the
targets of the generator; the nodes of 24 networks per level were resampled from half B;
the comparison uses the confidants of half B (up to five per respondent). Intervals come
from a bootstrap of half-B respondents within race, jointly with a bootstrap of the 24
networks ($B = 1000$). Table~\ref{tab:validation} reports observed and predicted values
and the relative error.

\paragraph{Confidants are not contacts.} A GSS respondent names about two people
(2.0 on average, 0 to 6) with whom they discuss important matters: the strongest ties of
a personal network that, in the model, has about 30 contacts. Two kinds of features are
therefore compared differently. Shares and distances (same race, sex and religion; age
and education differences) are compared with all the contacts of a node: the generator
assumes that confidants have the composition of the whole personal network, which is how
the targets enter. Strong ties tend to be more homophilous than weak ones
\citep{mcpherson2001}, so this assumption makes the homophily of the networks an upper
bound. Diversity within an ego network and closure among alters depend on which
contacts are compared. Each half-B respondent is matched to a node of the same race and
the same degree rank (the number of confidants as a proxy of network size, the
assumption behind the relative-degree target), and $k = \min(\text{confidants}, 5)$ of
its contacts are taken by three rules that bracket the unobserved selection by tie
strength: at random (the average contact); the $k$ nearest in the Gower distance of the
model (selection by similarity, the criterion of the influence gate); and the $k$ with
the most common neighbours with the ego (selection by overlap, the classic correlate of
tie strength; \citealp{granovetter1973}).

\paragraph{Homophily.} Out of sample, the targeted features hold: the share of
same-race confidants is 0.895 [0.874, 0.914] in half B and 0.896 in the networks, and
relative degree and same-race share fall within the interval for every group (the
intervals for Asian respondents are wide: there are 20 per half). The features that were
not targeted follow from the sex, religion, age and education coefficients of
\citet{smith2014}, estimated with the GSS 1985 and 2004 confidants, so they are not
independent of half B; what is tested is whether those coefficients, combined with the
race calibration and 30 contacts, give the confidant-level shares. Same religion
($-3\%$) and age difference ($+6\%$) are within the intervals; same sex is 5\% too low
and the education difference 9\% too high, small errors whose intervals exclude zero.
The racial diversity of random sets of $k$ contacts is 30\% higher than among
confidants: the networks spread cross-race ties evenly over egos, while in the survey
they concentrate in fewer ego networks. Diversity in religion and sex fits.

\paragraph{Closure.} Among the confidants of half B, 80.5\% of pairs know each other and
42.8\% are especially close. A tie of the model lies between those two definitions, so
we read them as a band. Under random selection no level comes close (8\% in x1, 16\% in
x2): confidants are not a random draw from the contacts. With the most embedded
contacts, x2 falls inside the band (63\%) and x1 at its lower edge (41\%); with the
nearest contacts, x2 is at the lower edge (38\%) and x1 below it (25\%). Closure also
varies with the size of the ego network. In half B it is flat to falling: egos with five
or more confidants minus egos with two or three, $-4.7$ points [$-11.7$, $+1.9$] for
knowing each other. x1 and x2 reproduce that profile under the three rules. A global
version of the rewiring, which doubles the total number of triangles instead of those of
every node, places the new triangles among high-degree nodes. Closure then rises with
ego-network size ($+16$ to $+28$ points, 17 to 33 points above the observed difference;
every interval excludes zero), and degree assortativity rises from 0.33 in x1 to 0.46
(0.41 in x2). This held-out profile is why x2 doubles the triangles of every node.

\paragraph{What the networks do not reproduce.} In the networks, the most embedded
contacts of a node are also much more alike in race and religion than real confidants
(Blau diversity 70--80\% lower). In the model, closure follows homophily; real
confidants who know each other are more varied. The level of closure itself is a
choice (twice that of x1), checked against these data but not estimated from them.

\begin{table}[htbp]
\centering
\scriptsize
\caption{Held-out validation: GSS 2004 confidants of half B vs.\ 24 networks per level
built with the targets of half A. Error: predicted/observed $- 1$; in bold when its 95\%
interval (bootstrap of respondents and networks) excludes zero. Rows with a rule in
parentheses compare $k = \min(\text{confidants}, 5)$ contacts of nodes matched by race
and degree rank; the other rows use all contacts. x2g: global rewiring (not used).}
\label{tab:validation}
% generated by scripts/09_network_validation.R; bold = 95% interval of the relative error excludes 0
\begin{tabular}{lcrrrrrr}
\hline
 & GSS 2004, half B & \multicolumn{2}{c}{x1} & \multicolumn{2}{c}{x2 (main)} & \multicolumn{2}{c}{x2g (global)} \\
Feature & observed [95\% CI] & pred. & error & pred. & error & pred. & error \\
\hline
Relative degree, White & 1.082 [1.047, 1.118] & 1.085 & +0\% & 1.085 & +0\% & 1.085 & +0\% \\
Relative degree, Black & 0.713 [0.568, 0.848] & 0.651 & -9\% & 0.651 & -9\% & 0.651 & -9\% \\
Relative degree, Hispanic & 0.868 [0.684, 1.064] & 0.923 & +6\% & 0.923 & +6\% & 0.923 & +6\% \\
Relative degree, Asian & 0.937 [0.593, 1.332] & 0.745 & -20\% & 0.745 & -20\% & 0.745 & -20\% \\
Same race, White & 0.948 [0.932, 0.965] & 0.950 & +0\% & 0.950 & +0\% & 0.950 & +0\% \\
Same race, Black & 0.804 [0.704, 0.890] & 0.873 & +9\% & 0.873 & +9\% & 0.873 & +9\% \\
Same race, Hispanic & 0.604 [0.480, 0.725] & 0.542 & -10\% & 0.542 & -10\% & 0.542 & -10\% \\
Same race, Asian & 0.528 [0.312, 0.757] & 0.615 & +17\% & 0.615 & +17\% & 0.615 & +17\% \\
Same race & 0.895 [0.874, 0.914] & 0.896 & +0\% & 0.896 & +0\% & 0.896 & +0\% \\
Same sex & 0.599 [0.572, 0.625] & 0.569 & \textbf{-5\%} & 0.569 & \textbf{-5\%} & 0.569 & \textbf{-5\%} \\
Same religion & 0.724 [0.690, 0.758] & 0.700 & -3\% & 0.700 & -3\% & 0.700 & -3\% \\
Age difference (years) & 11.138 [10.507, 11.846] & 11.643 & +5\% & 11.756 & +6\% & 11.757 & +6\% \\
Education difference (categories) & 1.268 [1.189, 1.353] & 1.372 & \textbf{+8\%} & 1.383 & \textbf{+9\%} & 1.382 & \textbf{+9\%} \\
Blau diversity, race (random) & 0.066 [0.051, 0.082] & 0.098 & \textbf{+49\%} & 0.086 & \textbf{+30\%} & 0.089 & \textbf{+36\%} \\
Blau diversity, race (nearest) & 0.066 [0.051, 0.082] & 0.004 & \textbf{-94\%} & 0.004 & \textbf{-93\%} & 0.004 & \textbf{-94\%} \\
Blau diversity, race (embedded) & 0.066 [0.051, 0.082] & 0.013 & \textbf{-80\%} & 0.020 & \textbf{-70\%} & 0.012 & \textbf{-82\%} \\
Blau diversity, religion (random) & 0.209 [0.184, 0.236] & 0.244 & \textbf{+17\%} & 0.223 & +7\% & 0.228 & +9\% \\
Blau diversity, religion (nearest) & 0.209 [0.184, 0.236] & 0.002 & \textbf{-99\%} & 0.004 & \textbf{-98\%} & 0.003 & \textbf{-98\%} \\
Blau diversity, religion (embedded) & 0.209 [0.184, 0.236] & 0.048 & \textbf{-77\%} & 0.040 & \textbf{-81\%} & 0.048 & \textbf{-77\%} \\
Blau diversity, sex (random) & 0.316 [0.296, 0.337] & 0.312 & -1\% & 0.306 & -3\% & 0.304 & -4\% \\
Blau diversity, sex (nearest) & 0.316 [0.296, 0.337] & 0.003 & \textbf{-99\%} & 0.003 & \textbf{-99\%} & 0.002 & \textbf{-99\%} \\
Blau diversity, sex (embedded) & 0.316 [0.296, 0.337] & 0.277 & \textbf{-12\%} & 0.263 & \textbf{-17\%} & 0.262 & \textbf{-17\%} \\
Pairs who know each other (random) & 0.805 [0.774, 0.833] & 0.084 & \textbf{-90\%} & 0.163 & \textbf{-80\%} & 0.187 & \textbf{-77\%} \\
Pairs who know each other (nearest) & 0.805 [0.774, 0.833] & 0.246 & \textbf{-70\%} & 0.380 & \textbf{-53\%} & 0.428 & \textbf{-47\%} \\
Pairs who know each other (embedded) & 0.805 [0.774, 0.833] & 0.413 & \textbf{-49\%} & 0.628 & \textbf{-22\%} & 0.735 & \textbf{-9\%} \\
Pairs especially close (random) & 0.428 [0.390, 0.464] & 0.084 & \textbf{-80\%} & 0.163 & \textbf{-62\%} & 0.187 & \textbf{-56\%} \\
Pairs especially close (nearest) & 0.428 [0.390, 0.464] & 0.246 & \textbf{-43\%} & 0.380 & \textbf{-11\%} & 0.428 & -0\% \\
Pairs especially close (embedded) & 0.428 [0.390, 0.464] & 0.413 & -3\% & 0.628 & \textbf{+47\%} & 0.735 & \textbf{+72\%} \\
\hline
 & \multicolumn{7}{l}{Closure, egos with 5+ confidants minus egos with 2--3 (pp): pred. \& pred.$-$obs.} \\
\hline
Pairs who know each other (random) & -4.7 [-11.7, +1.9] & -0.3 & +4.4 & -0.5 & +4.2 & +16.3 & \textbf{+21.0} \\
Pairs who know each other (nearest) & -4.7 [-11.7, +1.9] & -3.7 & +1.0 & -3.4 & +1.3 & +20.9 & \textbf{+25.6} \\
Pairs who know each other (embedded) & -4.7 [-11.7, +1.9] & -8.5 & -3.8 & -9.7 & -5.0 & +28.2 & \textbf{+32.9} \\
\hline
\end{tabular}

\end{table}
